\documentclass[11pt]{article}
\usepackage[margin=1in]{geometry}
\usepackage[T1]{fontenc}
\usepackage[utf8]{inputenc}
\usepackage{lmodern}
\usepackage{amsmath,amssymb}
\usepackage{booktabs}
\usepackage{graphicx}
\usepackage{xcolor}
\usepackage{hyperref}
\usepackage{float}

\hypersetup{
  colorlinks=true,
  linkcolor=blue!50!black,
  urlcolor=blue!50!black
}

\setlength{\parskip}{0.45em}
\setlength{\parindent}{0pt}

\title{\textbf{EVQ-Cosh: Core Theory and Experimental Story}}
\author{}
\date{March 6, 2026}

\begin{document}
\maketitle

\section*{Overview}

The central claim is simple: long-context extrapolation fails not only because
inference-time scaling is imperfect, but because standard geometric RoPE starts
from the wrong training-time frequency allocation. EVQ-Cosh reframes frequency
allocation as a variational inverse problem, derives a closed-form family, and
shows that geometric RoPE is the $\tau=0$ degenerate limit of that family.

The empirical picture is equally direct. In the cleanest extrapolation regimes,
EVQ outperforms learnable positional encodings despite using no extra
parameters. In the systems setting, the strongest result is not merely
``EVQ $>$ Geo'', but rather \textbf{EVQ + YaRN $\gg$ Geo + YaRN}. This suggests
that the hidden bottleneck for long-context systems is the training-time
frequency substrate, not only the inference-time scaling rule.

\section*{Core Thesis}

\textbf{Training-time frequency allocation is the hidden bottleneck for
long-context extrapolation.}

Standard geometric RoPE distributes frequencies uniformly in log space. That
choice is analytically convenient, but it is not the generic optimum of the
underlying interference problem. EVQ-Cosh instead reallocates channel density:
it compresses high-frequency spacing and expands low-frequency spacing. The
point is not to ``move frequencies to low frequencies''; mathematically that is
incorrect. The point is to reduce low-frequency phase collisions, which become
the structural bottleneck at long range.

This leads to three main consequences:

\begin{enumerate}
  \item Geometric RoPE should be treated as a special case, not as the default optimum.
  \item Extreme extrapolation quality can be improved by changing only the training-time frequency allocation.
  \item Inference-time scaling methods such as YaRN become much stronger when the underlying frequency allocation is repaired first.
\end{enumerate}

\section*{Theory in One Page}

\subsection*{1. Variational formulation}

Let $D(\Delta)$ denote a distance prior and let $\phi$ index the log-frequency
axis. The exact phase-interference kernel is
\[
K(\phi_1,\phi_2)
= \int D(\Delta)\cos(b^{-\phi_1}\Delta)\cos(b^{-\phi_2}\Delta)\,d\Delta.
\]

The only approximation step is a broadband surrogate:
\[
K(\phi_1,\phi_2) \approx \alpha\,\delta(\phi_1-\phi_2)+\beta\,\min(\phi_1,\phi_2).
\]

This turns the optimization over channel density $\rho(\phi)$ into the
variational objective
\[
J[\rho]
= \frac{\alpha}{2}\int \rho(\phi)^2\,d\phi
+ \frac{\beta}{2}\int\!\!\int \rho(\phi)\rho(\psi)\min(\phi,\psi)\,d\phi\,d\psi
- \mu\int \rho(\phi)b^{-2\phi}\,d\phi.
\]

\subsection*{2. Exact ODE and closed-form family}

The Euler--Lagrange equation reduces to
\[
\rho''(\phi)-\tau^2\rho(\phi)=\gamma b^{-2\phi},
\qquad \tau=\sqrt{\beta/\alpha}.
\]

Its solution is exact:
\[
\rho^*(\phi)=C_1\cosh(\tau\phi)+C_2\sinh(\tau\phi)+Pb^{-2\phi}.
\]

After CDF inversion, the EVQ-Cosh warp is
\[
\phi_k(\tau)
= 1-\frac{1}{\tau}\operatorname{arcsinh}\!\left((1-u_k)\sinh\tau\right).
\]

\subsection*{3. Geometric RoPE as a degenerate limit}

Taking $\tau\to 0$ yields
\[
\phi_k(\tau)\to u_k,
\]
which is exactly the uniform-quantile allocation of geometric RoPE. In this
sense, geometric RoPE is not a fundamentally distinct design; it is the zero
temperature limit of the EVQ family.

\subsection*{4. Waterbed interpretation}

The theory predicts a constrained tradeoff. EVQ reduces long-range low-frequency
collisions by spending positional budget away from high-frequency redundancy.
This is not a symmetric trade: high-frequency spacing is often partially
redundant, while low-frequency collisions are structurally destructive for very
long-range separation.

\subsection*{5. Scaling law}

The practical rule supported by the current evidence is
\[
\tau^*(L)\approx \frac{d_{\text{head}}}{\sqrt{L}}.
\]

In the current paper package, this scaling law should still be treated as an
empirical law or conjecture, not as a fully closed theorem. The exact ODE and
the geometric limit are the paper-grade theoretical anchors.

\section*{What The Experiments Establish}

\subsection*{A. Extreme extrapolation beats learnable PE}

In the 128-token DAPE-style PE-dominant regime, EVQ wins in the exact setting
where positional encoding quality is the dominant variable.

\begin{table}[H]
\centering
\small
\begin{tabular}{lccc}
\toprule
Method & Extra params & PPL@8K & vs Geo \\
\midrule
Geometric & 0 & 513.7 & --- \\
DAPE & 32 & 455.3 & -11.4\% \\
Learnable $\tau$ & 1 & 437.9 & -14.8\% \\
\textbf{EVQ $\tau=5.0$} & \textbf{0} & \textbf{333.7} & \textbf{-35.0\%} \\
\bottomrule
\end{tabular}
\caption{128-token training, 64$\times$ extrapolation to 8K. EVQ outperforms a learnable positional encoding family while using no additional parameters.}
\end{table}

This result matters because it isolates the positional encoding itself. It
shows that a closed-form frequency family can outperform a more flexible but
less structured learnable parameterization.

\subsection*{B. EVQ unlocks YaRN}

The main systems result is the interaction with YaRN. The strongest story is
not that EVQ is a slightly better standalone RoPE variant. The stronger story
is that EVQ changes how well inference-time scaling works.

\begin{table}[H]
\centering
\small
\begin{tabular}{lcccc}
\toprule
Method & PK@8K & PK@12K & PPL@8K & PPL@16K \\
\midrule
Geo & 41\% & 57\% & 161.9 & 253.2 \\
Geo + YaRN & 61\% & 59\% & 82.9 & 157.7 \\
EVQ & 53\% & 63\% & 150.3 & 229.5 \\
\textbf{EVQ + YaRN} & \textbf{100\%} & \textbf{79\%} & \textbf{70.9} & \textbf{107.5} \\
\bottomrule
\end{tabular}
\caption{Main systems comparison in the fair 10\% mix, scale=8 setting. EVQ+YaRN clearly dominates Geo+YaRN.}
\end{table}

The interpretation is straightforward. Geo compresses low-frequency spacing so
aggressively that inference-time rescaling still starts from a compromised
substrate. EVQ first repairs the low-frequency spacing, and YaRN can then act
on a much stronger base.

\subsection*{C. PE-dominant scaling-law confirmation}

At $L_{\text{train}}=256$, the scaling-law prediction points to $\tau^*=4.0$.
The experiment confirms this direction directly, and the YaRN interaction is
especially strong in this regime.

\begin{itemize}
  \item Raw: EVQ $\tau=4.0$ outperforms EVQ $\tau=2.0$ across all OOD lengths.
  \item With YaRN: EVQ4+YaRN reaches \textbf{99.6} PPL at $32\times$, versus \textbf{260.2} for Geo+YaRN.
  \item The implied gain is \textbf{-61.7\%} at the most extreme evaluation point.
\end{itemize}

\section*{Figures}

\begin{figure}[H]
\centering
\includegraphics[width=\textwidth]{../figs/fig2_evq_yarn_synergy.pdf}
\caption{Main systems figure. The strongest practical result is that EVQ unlocks YaRN: the training-time frequency allocation and inference-time scaling are complementary rather than redundant.}
\end{figure}

\begin{figure}[H]
\centering
\includegraphics[width=\textwidth]{../figs/fig3_pe_dominant_scaling.pdf}
\caption{PE-dominant and extreme extrapolation figure. The same theory explains both DAPE-style extrapolation and the stronger Phase 11 scaling-law regime.}
\end{figure}

\section*{Current Interpretation}

The evidence currently supports the following reading:

\begin{enumerate}
  \item The exact theoretical core is real: the ODE, its closed-form family, and the geometric limit are not heuristic curve fits.
  \item The strongest empirical message is systems-level: EVQ repairs the training-time positional substrate and thereby unlocks more of YaRN.
  \item The PE-dominant results are important because they show that this is not merely a downstream tuning trick; positional encoding quality alone can create large extrapolation gaps.
  \item Larger-scale continued pretraining remains supportive rather than primary, but it is directionally consistent: the long-range gain survives scale-up.
\end{enumerate}

\section*{What Remains Open}

Three open fronts are still worth closing:

\begin{itemize}
  \item a tighter perturbation-bound story for the broadband surrogate;
  \item a cleaner larger-scale text anchor on a real downstream benchmark;
  \item stronger cross-modal evidence if video is to raise the paper ceiling rather than remain supportive.
\end{itemize}

The current package is already enough to support a clear claim: long-context
extrapolation depends critically on how RoPE frequencies are allocated during
training, and repairing that allocation yields a closed-form method that is
both theoretically grounded and practically strong.

\end{document}
